\section{Hallucination des modèles multimodaux lors de l'extraction d'entités et de relations médicales}

Comme évoqué lors de la présentation des modèles multimodaux, l'utilisation des MLLMs dans les pipelines d'extraction d'information clinique introduit un risque spécifique : celui de générer des entités, des relations ou des valeurs plausibles mais absentes du document source. Ce phénomène, appelé hallucination, découle du fait que ces modèles interprètent le contenu à partir de leurs connaissances apprises, au lieu de se limiter à restituer le texte observé.

Cette section présente les principaux travaux consacrés à la compréhension, à la détection et à la réduction des hallucinations, en s'intéressant surtout à l'extraction structurée d'entités et de relations, où chaque information extraite doit être fidèle au document source pour pouvoir alimenter un graphe de connaissances.

\subsection{Définitions et taxonomies générales de l'hallucination}

La littérature distingue généralement deux grandes catégories d'hallucinations. Ji et al.~\cite{ji2023survey}, dans leur revue consacrée à la génération de texte, opposent les hallucinations intrinsèques, qui correspondent à des informations contredisant directement le contenu du document source, aux hallucinations extrinsèques, qui désignent des informations ne pouvant pas être vérifiées à partir du document, sans pour autant être nécessairement fausses. Huang et al.~\cite{huang2023survey} prolongent cette vision dans le cas des LLMs en distinguant les hallucinations de factualité, qui contredisent les connaissances du monde, des hallucinations de fidélité, qui s'écartent du contexte ou des instructions fournis au modèle. Ils identifient également les différentes sources d'hallucination tout au long du cycle de vie d'un modèle, depuis la collecte des données d'entraînement et la conception de l'architecture jusqu'à la phase d'inférence, et distinguent les méthodes visant à détecter les hallucinations de celles destinées à les atténuer. Dans le cadre de l'extraction d'information, c'est la fidélité qui importe : une information médicalement exacte mais absente du dossier du patient reste une erreur.

Dans le cas des modèles multimodaux, Bai et al.~\cite{bai2025survey} proposent une taxonomie plus spécifique adaptée aux interactions entre les modalités visuelle et textuelle. Leur étude recense les principales causes d'hallucination, liées aux données, au modèle, à l'entraînement ou à l'inférence, les benchmarks utilisés pour les évaluer ainsi que les différentes stratégies de mitigation. Les auteurs distinguent notamment les hallucinations portant sur la catégorie d'un objet, c'est-à-dire sur son existence même, sur ses attributs (valeur, quantité, couleur, position) ou encore sur les relations entre plusieurs objets présents dans une image ou un document scanné. Cette classification est particulièrement pertinente pour l'extraction d'informations à partir de documents médicaux numérisés, où une erreur sur une entité, une valeur clinique ou une relation peut se propager aux étapes suivantes du pipeline, notamment lors de la construction d'un graphe de connaissances.

Bien que ces taxonomies constituent une référence pour l'étude des hallucinations, elles ont principalement été élaborées pour des tâches de génération libre, telles que la description d'images, la réponse à des questions ou le résumé automatique. Elles prennent donc peu en compte les spécificités des tâches d'extraction structurée, dans lesquelles les informations produites doivent respecter un schéma prédéfini d'entités et de relations. Cette différence est importante, car les erreurs observées dans ces deux contextes ne sont pas de même nature. En extraction structurée, une hallucination ne se limite pas à l'invention d'une information absente du document ; elle peut également se traduire par une incohérence structurelle, par exemple lorsqu'une relation est créée entre deux entités incompatibles avec les contraintes du schéma ou lorsqu'une entité est associée à un type de relation qui ne lui est pas autorisé. Ce type d'erreur est critique dans la construction d'un graphe de connaissances, où la validité des relations dépend directement de la cohérence entre les types d'entités et les relations qui les relient.

\subsection{Hallucination dans l'extraction d'entités et de relations hors domaine médical}

Les travaux consacrés à la reconnaissance d'entités nommées (NER) et à l'extraction de relations (RE) à l'aide des LLMs évaluent généralement les performances à l'aide de métriques globales, telles que la précision, le rappel et le F1-score. En revanche, ils s'intéressent rarement à la nature des erreurs produites par les modèles et les analysent encore moins sous l'angle des hallucinations.

Tran~\cite{tran2025aspworkflow} propose un pipeline combinant un LLM et l'\textit{Answer Set Programming} (ASP), un paradigme de programmation déclarative fondé sur la logique, dans lequel les connaissances sont représentées sous forme de faits et de règles permettant de vérifier automatiquement le respect de contraintes prédéfinies. Dans cette approche, le LLM réalise l'extraction initiale des entités et des relations, tandis que le programme ASP applique un raisonnement symbolique afin de contrôler la cohérence des informations extraites, par exemple en vérifiant la compatibilité entre les types d'entités et les relations autorisées. Cette étape permet d'éliminer certaines incohérences structurelles avant la production du résultat final. Bien que ces erreurs ne soient pas explicitement qualifiées d'hallucinations par l'auteur, cette approche montre qu'une validation fondée sur des contraintes logiques peut renforcer la fiabilité des résultats produits par le modèle.

Sur le plan méthodologique, plusieurs approches de détection proposées dans la littérature peuvent être adaptées à la tâche d'extraction d'information. Certaines d'entre elles détectent les hallucinations au niveau des \textit{tokens} ou des \textit{spans}, c'est-à-dire de segments continus de texte correspondant à une entité ou à une expression d'intérêt, plutôt qu'à l'échelle de la phrase entière. Zhou et al.~\cite{zhou2021detecting} entraînent ainsi un modèle à prédire, pour chaque \textit{token} généré, s'il est appuyé ou non par le texte source. Mishra et al.~\cite{mishra2024fava} proposent une taxonomie fine des hallucinations (entité erronée, relation erronée, affirmation invérifiable, etc.) et un modèle, FAVA, qui repère les segments hallucinés puis les corrige à partir de documents récupérés. Dans une logique voisine, FActScore~\cite{min2023factscore} décompose un texte généré en faits atomiques, c'est-à-dire en affirmations élémentaires, et vérifie chacun d'eux de manière indépendante contre une source de connaissances. D'autres travaux exploitent les représentations internes des modèles de langage afin de repérer des comportements atypiques associés aux hallucinations. Azaria et Mitchell~\cite{azaria2023internal} montrent qu'un classifieur simple appliqué aux activations cachées d'un LLM permet de prédire si une affirmation produite est vraie ou fausse, et Chen et al.~\cite{chen2024inside} exploitent la cohérence des représentations internes entre plusieurs générations pour détecter les hallucinations. L'idée est que les représentations produites lors de la génération diffèrent lorsque le modèle s'appuie sur des informations présentes dans le document ou, au contraire, lorsqu'il génère un contenu non fondé sur celui-ci. Ces dernières approches supposent toutefois un accès aux états internes du modèle, ce qui n'est pas possible lorsque celui-ci est utilisé en boîte noire à travers une API. Enfin, Kadavath et al.~\cite{kadavath2022know} montrent que les LLMs sont en partie capables d'estimer eux-mêmes la probabilité que leur réponse soit correcte, mais cette auto-évaluation reste imparfaite et dépend de la tâche et de la formulation de la question.

Un travail plus récent, LATTE~\cite{zhang2026latte}, propose une nouvelle architecture visant à améliorer la stabilité de l'extraction conjointe des entités et des relations, notamment dans les phrases complexes. Les auteurs obtiennent des scores F1 de 93,8\,\% et 95,1\,\% sur les corpus NYT et WebNLG. Cette étude montre que les erreurs liées à l'identification des entités et des relations restent un défi important. Toutefois, leur approche repose sur une modification de l'architecture du modèle pendant son entraînement, ce qui n'est pas applicable à un modèle multimodal propriétaire utilisé en boîte noire.

Ces travaux mettent en évidence un manque : à notre connaissance, il n'existe pas encore de caractérisation systématique des hallucinations propres aux schémas d'extraction contraints, où l'erreur peut porter indépendamment sur la valeur d'une entité, son type, ou la cohérence structurelle du triplet entité--relation--entité qu'elle compose. La taxonomie de Mishra et al.~\cite{mishra2024fava}, qui distingue erreurs d'entité et erreurs de relation, s'en approche, mais elle porte sur du texte libre et ne tient pas compte des contraintes d'un schéma.

\subsection{Hallucination dans les MLLMs appliqués aux documents scannés}

Ce sous-domaine constitue le croisement le plus proche des configurations d'extraction directe par MLLM, sans passage par un OCR intermédiaire. Les premiers travaux sur l'hallucination des modèles vision-langage ont porté sur les images naturelles : Rohrbach et al.~\cite{rohrbach2018chair} ont introduit la métrique CHAIR, qui mesure la proportion d'objets mentionnés dans une légende mais absents de l'image, et Li et al.~\cite{li2023pope} ont proposé POPE, qui évalue l'hallucination d'objets en interrogeant le modèle par des questions fermées (\og y a-t-il un \ldots{} dans l'image ? \fg{}). Ces protocoles ne s'appliquent toutefois pas directement aux documents textuels, où l'unité d'information n'est pas un objet mais une valeur ou une mention écrite.

He et al.~\cite{he2025seeing} proposent KIE-HVQA, présenté comme le premier benchmark dédié à l'hallucination OCR dans la compréhension de documents dégradés. Il rassemble des cartes d'identité, des factures et des ordonnances, auxquelles sont appliquées des dégradations simulant des conditions réelles : flou de mouvement, faible contraste, occultation ou masquage d'une partie du texte. Les auteurs montrent que, face à ces dégradations, les MLLMs reconnaissent rarement leur propre incertitude. Au lieu de signaler qu'un caractère est illisible, ils s'appuient sur des biais linguistiques, c'est-à-dire sur les connaissances apprises par le modèle et sur les formulations les plus probables, et produisent des valeurs absentes de l'image. Pour réduire ce phénomène, ils entraînent Qwen2.5-VL-7B par apprentissage par renforcement avec une récompense qui distingue les caractères lisibles, partiellement masqués et illisibles, ces derniers devant être explicitement laissés vides. Le modèle obtenu dépasse GPT-4o d'environ 22 points de précision sans hallucination sur KIE-HVQA, tout en conservant ses performances OCR générales. Ce travail est particulièrement proche du contexte des comptes rendus scannés, car il inclut des documents médicaux (ordonnances) et montre que, pour un document dégradé, la bonne réponse est parfois de ne rien extraire. Il repose toutefois sur un réentraînement du modèle, ce qui n'est pas applicable à un modèle propriétaire utilisé en boîte noire.

Sur le plan mécanistique, Baek et al.~\cite{baek2025ocrheads} identifient au sein des modèles vision-langage des têtes d'attention spécialisées dans la lecture du texte présent dans l'image, qu'ils appellent têtes OCR. Ces têtes se distinguent des têtes de récupération (\textit{retrieval heads}) qui copient l'information dans un contexte textuel : elles se recouvrent peu avec elles, sont beaucoup moins rares, et leur activation est statique, c'est-à-dire directement liée à leur rôle de lecture, alors que les têtes de récupération sont activées de manière plus variable. Masquer une vingtaine de têtes OCR dans Qwen2-VL-7B fait chuter le F1 sur DocVQA de 76,3 à 59,7, ce qui confirme leur rôle central dans la lecture documentaire. Ces résultats offrent une piste de compréhension du comportement interne du modèle en complément des approches purement comportementales. Ils supposent cependant un accès aux poids et aux états internes du modèle, et les auteurs ne les évaluent pas directement pour la réduction des hallucinations.

Le travail le plus proche de la configuration visée par une extraction clinique directe est celui de RAPTOR+~\cite{raptorplus2026}, qui cible spécifiquement le traitement de documents cliniques scannés de mauvaise qualité (tampons administratifs, ratures, mentions surchargées, modèles hétérogènes d'un centre de référencement à l'autre). Les auteurs y proposent un cadre d'évaluation quantifiant conjointement la précision de l'extraction structurée et la qualité de localisation de la preuve visuelle, et caractérisent deux modes d'échec cliniquement critiques : les extractions non ancrées et la mauvaise localisation de l'évidence.

Le principe de l'ancrage visuel, qui consiste à appuyer les réponses du modèle sur des éléments réellement présents dans l'image, a été mis en œuvre par Woodpecker~\cite{woodpecker2023} pour corriger les hallucinations d'un MLLM. Cette approche utilise des boîtes englobantes (\textit{bounding boxes}), c'est-à-dire des rectangles qui délimitent la position des objets ou des régions d'intérêt dans une image, comme preuves visuelles. Elle s'appuie sur un pipeline en cinq étapes : l'extraction de concepts clés, la formulation de questions, la validation des connaissances visuelles par des modèles experts, la génération d'affirmations visuelles, puis la correction des hallucinations. Woodpecker présente l'intérêt de ne nécessiter aucun réentraînement : la correction intervient après la génération, ce qui la rend applicable à un modèle utilisé en boîte noire. Dans le domaine médical, MedSIGHT~\cite{chang2026medsight} reprend ce principe en générant des boîtes englobantes à partir de masques de segmentation, mais son application reste limitée à l'imagerie radiologique.

L'analyse de ces travaux met en évidence une limite : aucun ne combine à la fois un document clinique textuel scanné et dégradé, une extraction structurée des entités et des relations, ainsi qu'un mécanisme d'ancrage dans la source. Parmi les approches existantes, RAPTOR+~\cite{raptorplus2026} est la plus proche de cette configuration. Toutefois, elle ne produit pas de graphe de connaissances et n'intègre ni validation symbolique ni validation ontologique.

\subsection{Hallucination dans l'extraction d'information clinique à partir de texte propre}

La revue systématique de Fraile Navarro et al.~\cite{fraile2023review}, portant sur 94 études consacrées à la reconnaissance d'entités et à l'extraction de relations dans les textes cliniques, montre que la majorité des travaux reposent sur l'apprentissage automatique, que les approches à base de règles restent minoritaires, et que seules trois études rapportent une utilisation en dehors d'un cadre expérimental. Avec l'arrivée des LLMs, Agrawal et al.~\cite{agrawal2022fewshot} ont montré que ces modèles sont de bons extracteurs d'information clinique en apprentissage à partir de quelques exemples (\textit{few-shot}), sans données annotées spécifiques, à condition de contrôler la forme et le contenu de leurs sorties.

Ce contrôle est d'autant plus nécessaire que les LLMs hallucinent aussi dans le domaine médical. Pal et al.~\cite{pal2023medhalt} le mettent en évidence avec Med-HALT, un benchmark qui évalue les hallucinations des LLMs sur des tâches de raisonnement médical et de restitution de connaissances. Wang et al.~\cite{wang2024feature} montrent, à partir d'un pipeline modulaire d'extraction de caractéristiques cliniques appliqué à 25\,709 dossiers de grossesse, que les \textit{feature hallucinations}, c'est-à-dire la génération de caractéristiques cliniques incorrectes ou inexistantes, représentent un enjeu important pour la sécurité des systèmes. De son côté, CLINES~\cite{yang2025clines} met en évidence que l'association des entités extraites aux terminologies médicales UMLS (\textit{Unified Medical Language System}) et SNOMED CT (\textit{Systematized Nomenclature of Medicine -- Clinical Terms}), deux référentiels utilisés pour normaliser les concepts cliniques, reste souvent imprécise. Les auteurs observent notamment des codes hallucinés et soulignent que les limites de longueur du contexte compliquent le raisonnement sur les notes cliniques volumineuses.

Dans une tâche proche, celle de la génération de résumés cliniques, une étude basée sur l'atomisation des faits, qui consiste à décomposer un texte en faits élémentaires avant leur traitement, montre qu'extraire les faits avant de les recombiner peut, de manière inattendue, augmenter les hallucinations et les omissions au lieu de les réduire~\cite{bn2025fact}. Enfin, un système de vérification factuelle ne faisant pas appel à un LLM lors de l'étape d'évaluation~\cite{zeba2026mitigating} montre qu'un mécanisme déterministe, c'est-à-dire reposant sur des règles produisant toujours le même résultat pour une même entrée, combinant plusieurs vérifications de cohérence logique, peut constituer une alternative intéressante.

\subsection{Auto-correction et correction après génération}

Une autre famille de travaux cherche non plus seulement à détecter les hallucinations, mais à les corriger après la génération. Madaan et al.~\cite{madaan2023selfrefine} proposent Self-Refine, dans lequel le modèle critique sa propre réponse puis la réécrit de manière itérative. Dhuliawala et al.~\cite{dhuliawala2024cove} proposent Chain-of-Verification : le modèle génère d'abord une réponse, puis formule des questions de vérification, y répond indépendamment et révise sa réponse initiale en conséquence.

La portée de ces approches est cependant discutée. Huang et al.~\cite{huang2024cannot} montrent qu'un LLM ne parvient pas à améliorer son raisonnement par auto-correction lorsqu'il ne dispose d'aucun retour extérieur, et que ses performances peuvent même se dégrader. La revue critique de Kamoi et al.~\cite{kamoi2024selfcorrection} conclut que l'auto-correction n'est fiable que lorsque le modèle reçoit un retour provenant d'une source extérieure fiable, par exemple un outil, un exécuteur de code ou un document de référence, ou après un affinage spécifique à grande échelle. Ces résultats sont importants pour l'extraction d'information : demander au modèle de vérifier lui-même ses extractions ne garantit pas leur fidélité au document, et la vérification gagne à s'appuyer sur des signaux extérieurs au modèle, tels que le texte source ou les contraintes d'un schéma.

\subsection{Ontologies, contraintes et réparation de graphes comme mécanisme de validation}

La vérification de la cohérence structurelle d'un graphe est un domaine bien établi en représentation des connaissances. La recommandation SHACL (\textit{Shapes Constraint Language}) du W3C~\cite{shacl2017} permet de déclarer les contraintes qu'un graphe RDF doit respecter, telles que les types autorisés, les propriétés attendues, les cardinalités ou le domaine et l'image d'une relation, et produit un rapport listant chaque violation détectée. Ce type de validation détecte les incohérences, mais ne les répare pas. En outre, un graphe peut respecter toutes les contraintes déclarées tout en contenant une information absente du document : la validité structurelle ne garantit pas la fidélité à la source.

La réparation des erreurs a été étudiée principalement dans le domaine des bases de données et des graphes de connaissances. Paulheim~\cite{paulheim2017refinement} recense les méthodes de raffinement de graphes de connaissances, qu'il répartit en deux familles : la complétion, qui ajoute les faits ou les types manquants, et la détection et la correction d'erreurs, qui retypent ou retirent les éléments faux. Il souligne également la difficulté d'évaluer ces méthodes en l'absence de vérité de terrain complète. En bases de données, la réparation consiste à rechercher la modification minimale qui fait disparaître les violations de contraintes. Chu et al.~\cite{chu2013holistic} traitent l'ensemble des violations simultanément, puisqu'une même valeur erronée peut en provoquer plusieurs, Fan et al.~\cite{fan2016gfd} définissent des dépendances fonctionnelles adaptées aux graphes, et HoloClean~\cite{rekatsinas2017holoclean} combine contraintes, statistiques et inférence probabiliste pour proposer des réparations. Ces méthodes ont toutefois été conçues pour des données déjà établies, où les erreurs sont rares et isolées, et elles ne s'appuient pas sur un document source pour justifier une correction.

Plusieurs travaux récents mobilisent par ailleurs les ontologies biomédicales comme mécanisme de validation indépendant du modèle générateur. Par exemple, un pipeline combinant une extraction guidée, la génération d'un schéma RDF/OWL et un consensus multi-LLM, où plusieurs modèles comparent leurs résultats afin d'obtenir une décision plus fiable, attribue des rôles différents à chaque modèle pour améliorer la robustesse de l'extraction~\cite{das2026clinical}. Ali et al.~\cite{ali2026ontology} construisent une ontologie RDF/OWL et un graphe de connaissances à partir de données hospitalières, puis les intègrent à un LLM pour ancrer ses réponses dans une tâche de question-réponse clinique. Sur 60 questions, l'exactitude passe de 37\,\% (ChatGPT-4) et 52\,\% (DeepSeek-R1) à 98\,\%, et le taux d'hallucination de 63\,\% et 48\,\% à 1,7\,\%. Ces résultats restent préliminaires, compte tenu de la taille réduite de l'évaluation. Ce travail porte en outre sur la génération de réponses, et non sur la validation d'un graphe extrait de documents.

En dehors du domaine médical, CyberBOT~\cite{zhao2025cyberbot} applique une étape de vérification après la génération des réponses en s'appuyant sur un raisonnement en logique de description, un formalisme utilisé pour représenter et vérifier automatiquement les connaissances d'une ontologie. Cette approche améliore de manière significative la fidélité des résultats. De même, une étude combinant une décomposition structurée du problème et des raisonneurs logiques, c'est-à-dire des outils capables de vérifier automatiquement la cohérence des connaissances, montre des gains significatifs par rapport au \textit{prompting} standard~\cite{sadowski2026structured}.

Une étude contrôlée sur l'ancrage par graphe de connaissances, qui consiste à utiliser un graphe de connaissances pour vérifier ou enrichir les réponses du modèle, dans une tâche de question-réponse clinique~\cite{mandarapu2026knowledge} montre que cette approche n'améliore les performances que lorsque les informations recherchées sont hors distribution, c'est-à-dire absentes ou peu représentées dans les données d'entraînement du modèle. En revanche, elle peut parfois dégrader les performances d'un modèle déjà performant.

\subsection{Hybrid-Code v2 : génération neuronale et validation symbolique indépendante}

Un travail neuro-symbolique récent, Hybrid-Code v2~\cite{hybrid2512v2}, combine un modèle neuronal avec un mécanisme de validation symbolique indépendant. Ce système génère d'abord des candidats de codes ICD-10, une classification internationale des maladies utilisée pour coder les diagnostics médicaux, puis les vérifie à l'aide d'une base de connaissances symbolique. La validation repose sur plusieurs niveaux de contrôle, notamment le format du code, son ancrage dans le document, la détection des négations, la cohérence temporelle et des règles d'exclusion. Les auteurs montrent que cette approche élimine complètement les erreurs de format (hallucinations de type I) et réduit les erreurs cliniquement incorrectes mais syntaxiquement valides (hallucinations de type II), sans toutefois les supprimer totalement.

Ce travail montre qu'il est possible de combiner une génération neuronale avec une validation symbolique indépendante. Toutefois, il s'intéresse à une tâche de classification fermée, c'est-à-dire au choix d'un code parmi une nomenclature existante de plus de 70\,000 entrées, à partir d'un texte clinique déjà propre. Contrairement aux tâches d'extraction de connaissances, il ne nécessite ni l'identification des entités et des relations ni la construction d'un graphe de connaissances.

\subsection{Architectures multi-agents pour la détection d'hallucination clinique}

L'idée de faire collaborer plusieurs modèles pour améliorer la fiabilité des réponses a d'abord été étudiée hors du domaine médical. Du et al.~\cite{du2024debate} montrent que faire débattre plusieurs instances d'un LLM, chacune critiquant les réponses des autres sur plusieurs tours, améliore la factualité et le raisonnement par rapport à un modèle seul.

Dans le domaine clinique, CuraView~\cite{ye2026curaview}, Dialectic-Med~\cite{lu2026dialectic} et BLUEmed~\cite{you2026bluemed} proposent des architectures multi-agents, dans lesquelles plusieurs agents collaborent en assurant des rôles complémentaires, comme l'extraction, la vérification et la correction. Ces approches reposent sur un mécanisme de génération suivi d'une vérification croisée des résultats. Toutefois, elles n'organisent pas encore les décisions des agents selon une typologie des erreurs ni selon leur origine, qu'il s'agisse d'informations propres au patient ou de connaissances médicales générales. De plus, lorsque les agents vérificateurs sont eux-mêmes des modèles de langage, la vérification hérite du risque qu'elle cherche à contrôler : un agent juge peut lui aussi halluciner. Ce constat rejoint les conclusions de Kamoi et al.~\cite{kamoi2024selfcorrection} sur la nécessité d'un retour extérieur fiable.

\subsection{Vérification sans modèle juge : cohérence, contexte clinique et décision sélective}

Face à ce risque de circularité, une autre famille de travaux, plus ancienne pour certains, repose sur des signaux explicites et reproductibles : la cohérence entre plusieurs générations, la détection de règles linguistiques dans le texte clinique, la combinaison pondérée d'heuristiques et la possibilité pour un système de s'abstenir. Ces travaux ne ciblent pas tous directement l'hallucination des MLLMs, mais ils fournissent des briques méthodologiques transposables à l'extraction structurée.

\paragraph{Cohérence entre générations.}
Wang et al.~\cite{wang2023selfconsistency} proposent la \textit{self-consistency} : au lieu de conserver une seule réponse, le modèle génère plusieurs chemins de raisonnement, puis la réponse la plus fréquente est retenue par vote majoritaire. Les auteurs montrent que cette stratégie améliore les performances sur des tâches de raisonnement arithmétique et de sens commun. Manakul et al.~\cite{manakul2023selfcheckgpt} exploitent une intuition voisine pour la détection d'hallucinations avec SelfCheckGPT : si le modèle connaît réellement une information, plusieurs générations stochastiques tendent à être cohérentes entre elles, alors qu'une information hallucinée tend à varier d'une génération à l'autre. Farquhar et al.~\cite{farquhar2024semantic} affinent cette idée avec l'entropie sémantique, qui regroupe les générations ayant le même sens avant de mesurer leur dispersion, afin de ne pas confondre une simple variation de formulation avec une incertitude réelle. Ces approches fonctionnent sans base de connaissances externe et, pour SelfCheckGPT, en boîte noire, ce qui les rend compatibles avec un modèle propriétaire accessible uniquement par API, au prix de plusieurs appels au modèle. Toutefois, la cohérence ne garantit pas la fidélité au document : un modèle peut produire de manière répétée la même information plausible mais absente de la source. Dans un contexte d'extraction, le vote doit donc être combiné à une vérification d'ancrage dans le document.

\paragraph{Négation, hypothèse, temporalité et sujet de l'information.}
Chapman et al.~\cite{chapman2001negex} proposent NegEx, un algorithme fondé sur des expressions régulières qui détermine si un concept clinique est nié à partir de déclencheurs (par exemple \og no \fg{} ou \og without \fg{}) situés dans une fenêtre de mots autour de la mention, initialement sur des comptes rendus de sortie. Harkema et al.~\cite{harkema2009context} étendent ce principe avec ConText, qui détermine trois propriétés contextuelles d'une mention : la négation, le statut temporel (récent, historique ou hypothétique) et le sujet de l'information (\textit{experiencer}), c'est-à-dire le fait que le constat concerne le patient ou une autre personne. Ces algorithmes sont déterministes et interprétables. Leur intérêt pour l'étude des hallucinations tient au fait qu'ils permettent de détecter un type d'erreur invisible à une simple vérification de présence dans le texte : un modèle qui extrait \og fièvre \fg{} à partir de \og absence de fièvre \fg{} produit une entité littéralement présente dans le document, mais un fait clinique faux. Ces méthodes ont été conçues pour l'anglais ; leur application au français suppose l'adaptation des listes de déclencheurs, comme l'ont fait Deléger et Grouin~\cite{deleger2012negation} pour la détection de la négation dans des notes cliniques françaises.

\paragraph{Combinaison pondérée d'heuristiques.}
Ratner et al.~\cite{ratner2017snorkel} proposent Snorkel, un système de supervision faible dans lequel des experts écrivent des fonctions d'étiquetage (\textit{labeling functions}), c'est-à-dire des règles heuristiques bruitées, dont les sorties sont combinées par un modèle qui estime la fiabilité de chacune afin de produire des étiquettes d'entraînement. L'intérêt de ce travail, au-delà de la création de données d'entraînement, réside dans l'idée qu'un ensemble de règles individuellement imparfaites peut être agrégé en une décision plus fiable, à condition d'estimer le poids de chaque règle plutôt que de les traiter toutes de manière identique.

\paragraph{Décision sélective.}
L'idée de permettre à un système de refuser de décider est ancienne : Chow~\cite{chow1970reject} formalise dès 1970 le compromis entre le taux d'erreur et le taux de rejet d'un classifieur. Geifman et El-Yaniv~\cite{geifman2017selective} transposent ce principe aux réseaux de neurones profonds avec la classification sélective : le système peut s'abstenir lorsque sa confiance est insuffisante, et un seuil sur une fonction de confiance est choisi de façon à garantir, avec une forte probabilité, un niveau de risque fixé à l'avance, tout en conservant la plus grande couverture possible. Ce compromis entre risque et couverture est bien adapté au domaine médical, où il est souvent préférable de ne pas produire une information plutôt que de produire une information fausse. Transposé à l'extraction, il correspond à un compromis explicite entre précision et rappel.

\paragraph{Bilan.}
Ces quatre familles de travaux couvrent chacune une partie du problème : la cohérence entre générations cible l'instabilité du modèle, les algorithmes de type NegEx et ConText ciblent les erreurs de contexte et de temporalité, la supervision faible fournit un cadre pour combiner des indices hétérogènes, et la décision sélective permet de contrôler explicitement le niveau de risque.

\subsection{Benchmarks et métriques d'évaluation}

Plusieurs benchmarks, c'est-à-dire des jeux de données et des protocoles de référence utilisés pour évaluer les modèles, sont aujourd'hui consacrés à l'évaluation des hallucinations en contexte clinique. Med-HALT~\cite{pal2023medhalt} évalue les LLMs sur des tests de raisonnement médical et de restitution de connaissances issus d'examens de médecine. ClinHallu~\cite{yang2026clinhallu} identifie l'étape du raisonnement où apparaît l'hallucination (perception, connaissance ou intégration) au lieu d'évaluer uniquement la réponse finale. Pour les modèles multimodaux, CHAIR~\cite{rohrbach2018chair} et POPE~\cite{li2023pope} mesurent l'hallucination d'objets dans des images naturelles, et KIE-HVQA~\cite{he2025seeing} l'hallucination de lecture sur des documents dégradés. Toutefois, aucun n'est spécifiquement conçu pour l'extraction d'entités et de relations à partir de documents cliniques scannés présentant de fortes contraintes structurelles.

Concernant les métriques d'évaluation, la plupart des travaux d'extraction se limitent à la précision, au rappel et au F1-score. Or, la précision mesure indirectement le taux d'hallucination, puisque toute information extraite absente de la référence est comptée comme un faux positif, mais elle ne distingue pas une information inventée d'une information correcte mal typée ou mal délimitée. Des métriques plus fines ont été proposées : FActScore~\cite{min2023factscore} mesure la proportion de faits atomiques appuyés par une source, et Ghanem et Cruz~\cite{ghanem2024enhancing} proposent un cadre reposant sur une mesure de similarité entre graphes basée sur BERTScore, une métrique qui compare la similarité sémantique entre deux textes à l'aide de représentations produites par BERT. Les auteurs recommandent un seuil de 95\,\% pour considérer que deux graphes correspondent, dans le cadre de la réduction des hallucinations et des omissions lors de la construction de graphes de connaissances à partir de textes génériques. Bien que cette étude ne porte pas sur le domaine clinique, elle constitue une référence importante pour l'évaluation de la fidélité structurelle, au-delà des métriques classiques telles que la précision et le rappel.

\subsection{Analyse critique et synthèse}

L'analyse de la littérature fait apparaître une couverture inégale des différentes dimensions du problème. L'ancrage visuel sur documents dégradés est partiellement couvert, principalement par RAPTOR+~\cite{raptorplus2026}, mais hors du cadre d'une extraction structurée vers un graphe de connaissances. La validation ontologique est bien établie sur texte propre, et la réparation sous contraintes est bien étudiée en bases de données, mais sans exiger que la correction soit justifiée par un document source. Les méthodes d'auto-correction et les architectures multi-agents introduisent des mécanismes de vérification, mais sans typologie fine des erreurs et en s'appuyant le plus souvent sur des modèles de langage comme juges, dont la fiabilité dépend de l'existence d'un retour extérieur~\cite{kamoi2024selfcorrection}. Hybrid-Code v2~\cite{hybrid2512v2} atteint une garantie formelle sur une catégorie d'erreur, mais sur une tâche de classification fermée qui ne pose jamais la question de la frontière entre entité et relation. Enfin, les méthodes de vérification sans modèle juge restent étudiées isolément, rarement au sein d'un pipeline de construction de graphe de connaissances à partir d'un MLLM.

Cette cartographie met ainsi en évidence un espace de recherche encore ouvert à l'intersection de trois éléments rarement combinés dans la littérature actuelle : le traitement direct de documents cliniques scannés, l'extraction structurée d'entités et de relations vers un graphe de connaissances, et une validation combinant ancrage dans la source, cohérence symbolique et vérification explicite de la stabilité de la frontière entité--relation. LATTE~\cite{zhang2026latte} aborde bien cette instabilité, mais au niveau de l'architecture du modèle et non en validation \textit{post-hoc}.


\subsection{Synthèse des hallucinations identifiées dans la littérature}

Le Tableau~\ref{tab:hallucinations} synthétise les principaux types d'hallucinations identifiés dans cette section, les travaux qui les étudient et les approches proposées pour les réduire. Il montre que les hallucinations ne se limitent pas à l'invention d'informations : elles peuvent affecter la perception visuelle, les entités, les relations, le contexte clinique de la mention, le raisonnement ou la cohérence structurelle du graphe.


\begin{table*}[ht]
\centering
\footnotesize
\renewcommand{\arraystretch}{1.35}
\begin{tabular}{|p{2.8cm}|p{4.4cm}|p{3.5cm}|p{4.3cm}|}
\hline
\textbf{Type d'hallucination} &
\textbf{Description} &
\textbf{Travaux} &
\textbf{Approches proposées} \\
\hline

Hallucination visuelle et de lecture &
Le modèle génère des informations absentes du document ou interprète incorrectement les éléments visuels, en particulier lorsque la qualité de l'image est dégradée. &
CHAIR \cite{rohrbach2018chair}, POPE \cite{li2023pope}, KIE-HVQA \cite{he2025seeing}, Woodpecker \cite{woodpecker2023}, MedSIGHT \cite{chang2026medsight}, RAPTOR+ \cite{raptorplus2026} &
Ancrage visuel, boîtes englobantes (\textit{bounding boxes}), segmentation, abstention sur les caractères illisibles et vérification après génération.\\
\hline

Hallucination d'entité &
Erreur sur la valeur d'une entité ou sur son type. &
LATTE \cite{zhang2026latte}, Wang et al. \cite{wang2024feature}, FAVA \cite{mishra2024fava} &
Amélioration de l'extraction, stabilisation des frontières de décision, détection au niveau des segments, validation des entités extraites.\\
\hline

Hallucination de relation &
Relation incorrecte entre deux entités ou mauvaise association entre les concepts extraits. &
LATTE \cite{zhang2026latte}, FAVA \cite{mishra2024fava} &
Extraction conjointe entité--relation, amélioration de la stabilité des frontières et détection des relations erronées.\\
\hline

Hallucination contextuelle &
Fait extrait alors qu'il est nié, hypothétique ou qu'il ne concerne pas le patient. &
NegEx \cite{chapman2001negex}, ConText \cite{harkema2009context}, Deléger et Grouin \cite{deleger2012negation}, Hybrid-Code v2 \cite{hybrid2512v2} &
Détection de déclencheurs de négation et d'hypothèse, identification du sujet de l'information.\\
\hline

Hallucination temporelle &
Information historique (antécédent) interprétée comme un fait actuel. &
ConText \cite{harkema2009context}, Hybrid-Code v2 \cite{hybrid2512v2} &
Détermination du statut temporel de la mention, vérification de la cohérence temporelle.\\
\hline

Hallucination ontologique &
Association erronée entre une entité extraite et un concept médical normalisé (UMLS, SNOMED CT ou ICD-10). &
CLINES \cite{yang2025clines}, Das et al. \cite{das2026clinical}, Hybrid-Code v2 \cite{hybrid2512v2} &
Validation ontologique, bases de connaissances symboliques, consensus multi-LLM et vérification des codes.\\
\hline

Hallucination de raisonnement &
Erreur apparaissant lors du raisonnement clinique ou de l'intégration des informations. &
Med-HALT \cite{pal2023medhalt}, ClinHallu \cite{yang2026clinhallu}, Dialectic-Med \cite{lu2026dialectic}, Zeba et al. \cite{zeba2026mitigating} &
Raisonnement multi-agents, vérification logique et validation déterministe.\\
\hline

Hallucination factuelle &
Production de faits médicaux inexacts ou omission d'informations pertinentes malgré un texte cohérent. &
BN et al. \cite{bn2025fact}, Wang et al. \cite{wang2024feature}, FActScore \cite{min2023factscore}, SelfCheckGPT \cite{manakul2023selfcheckgpt}, entropie sémantique \cite{farquhar2024semantic} &
Atomisation des faits, extraction modulaire, vérification factuelle et cohérence entre générations.\\
\hline

Hallucination structurelle &
Incohérences dans le graphe de connaissances, relations incompatibles ou non-respect des contraintes du schéma ou de l'ontologie. &
SHACL \cite{shacl2017}, Paulheim \cite{paulheim2017refinement}, HoloClean \cite{rekatsinas2017holoclean}, CyberBOT \cite{zhao2025cyberbot}, Sadowski et al. \cite{sadowski2026structured}, Ghanem et Cruz \cite{ghanem2024enhancing}, Hybrid-Code v2 \cite{hybrid2512v2} &
Validation par contraintes, réparation minimale, raisonnement en logique de description, raisonneurs logiques et mesure de similarité entre graphes.\\
\hline

\end{tabular}
\caption{Synthèse des principaux types d'hallucinations identifiés dans la littérature et des approches proposées pour les réduire.}
\label{tab:hallucinations}
\end{table*}

\section{Conclusion du chapitre}

Ce chapitre a montré que les briques nécessaires à une extraction clinique fiable existent séparément dans la littérature : extraction multimodale, annotation guidée par des règles explicites, ancrage dans la source, validation par contraintes et réparation de graphes, détection du contexte clinique, vote entre générations et décision sélective. Il a également montré que l'auto-correction par le modèle lui-même ou par d'autres modèles juges reste fragile en l'absence d'un retour extérieur fiable. En revanche, ces briques ont rarement été assemblées dans un même pipeline appliqué à des documents cliniques scannés et orienté vers la construction d'un graphe de connaissances. Le chapitre suivant présente l'approche proposée pour répondre à cette limite.
